% !TeX program = xelatex
\documentclass[11pt,a4paper]{ctexart}

\usepackage[left=2.4cm,right=2.4cm,top=2.1cm,bottom=2.1cm]{geometry}
\usepackage{amsmath,amssymb,mathtools,braket}
\usepackage{xcolor}
\usepackage[inline]{enumitem}
\usepackage{fancyhdr}

% ============================================================
% 基本信息：只需修改下面几项
% ============================================================
\newcommand{\coursename}{高等量子力学}
\newcommand{\assignmentnumber}{第 1 次作业}
\newcommand{\teachername}{\underline{\hspace{2.8cm}}}
\newcommand{\studentname}{\underline{\hspace{2.8cm}}}
\newcommand{\studentid}{\underline{\hspace{3.0cm}}}

% ============================================================
% 数学公式书写规范
% ============================================================
% 1. 行内公式使用一对美元符号。
% 2. 独立成行的公式统一使用 equation 环境，不使用反斜杠方括号形式。
% 3. 虚数单位 i、自然常数 e 和微分符号 d 均使用直立体 \mathrm。

% ============================================================
% 页面样式
% ============================================================
\setlength{\parindent}{2em}
\setlength{\parskip}{0.25em}
\setlength{\headheight}{14pt}
\pagestyle{fancy}
\fancyhf{}
\lhead{\coursename}
\rhead{\assignmentnumber}
\cfoot{\thepage}
\renewcommand{\headrulewidth}{0.4pt}

% ============================================================
% 列表参数
% ============================================================
\setlist[enumerate,1]{
  label=(\alph*),
  leftmargin=3.2em,
  itemsep=0.55em,
  topsep=0.45em
}

\newlist{inlineenum}{enumerate*}{1}
\setlist[inlineenum,1]{
  label=(\alph*),
  labelsep=0.5em,
  itemjoin={；\quad},
  itemjoin*={；\quad}
}

% ============================================================
% 题目与回答环境
% ============================================================
\newcounter{problem}
\newenvironment{problem}
  {\refstepcounter{problem}%
   \par\medskip\noindent
   \begingroup
   \setlength{\fboxsep}{3pt}%
   \setlength{\fboxrule}{0.8pt}%
   \fbox{\bfseries 题目\chinese{problem}}%
   \endgroup
   \quad\ignorespaces}
  {\par\medskip}

\newenvironment{answer}
  {\par\smallskip\noindent
   \begingroup
   \setlength{\fboxsep}{4pt}%
   \colorbox{black}{\textcolor{white}{\bfseries 回答}}%
   \endgroup
   \quad\ignorespaces}
  {\par\smallskip\noindent\rule{\linewidth}{0.4pt}\medskip}

\begin{document}

\begin{center}
  {\LARGE\bfseries \coursename\quad \assignmentnumber}\par
  \vspace{0.7em}
  \small
  教师：\teachername\qquad
  姓名：\studentname\qquad
  学号：\studentid
\end{center}
\vspace{0.4em}

\begin{problem}
设三维希尔伯特空间基矢为 $\lvert a_1\rangle$、$\lvert a_2\rangle$、$\lvert a_3\rangle$。
\begin{enumerate}
  \item 写出 $\lvert a_1\rangle$、$\lvert a_3\rangle\langle a_2\rvert$ 的矩阵表示。
  \item 求
  \begin{equation}
    \lvert a_3\rangle\langle a_2\rvert
    +\lvert a_1\rangle\langle a_1\rvert
    +\lvert a_2\rangle\langle a_3\rvert
  \end{equation}
  的本征值和本征矢。
\end{enumerate}
\end{problem}

\begin{answer}
% 在此处作答
\end{answer}

\begin{problem}
根据厄米共轭算符的定义，证明：
\begin{enumerate}
  \item $\bigl(\lvert\psi\rangle\langle\phi\rvert\bigr)^\dagger
        =\lvert\phi\rangle\langle\psi\rvert$；
  \item $\bigl(\hat A^\dagger\bigr)^\dagger=\hat A$；
  \item $\bigl(\hat A\hat B\bigr)^\dagger=\hat B^\dagger\hat A^\dagger$。
\end{enumerate}
\end{problem}

\begin{answer}
% 在此处作答
\end{answer}

\begin{problem}
导出坐标表象下，下列算符矩阵元的表达式：
\begin{enumerate}
  \item $(\hat x\hat p)^2$；
  \item $\mathrm{e}^{\mathrm{i}a\hat p}$；
  \item $\mathrm{e}^{\mathrm{i}p\hat x}$。
\end{enumerate}
\end{problem}

\begin{answer}
% 在此处作答
\end{answer}

\begin{problem}
证明：
\begin{enumerate}
  \item 厄米算符的本征值为实数；
  \item 反厄米算符的本征值为纯虚数；
  \item 幺正算符的本征值为模为 $1$ 的复数。
\end{enumerate}
\end{problem}

\begin{answer}
% 在此处作答
\end{answer}

\begin{problem}
设
\begin{equation}
  F_{x,x'}=-\mathrm{i}x^2\frac{\partial}{\partial x}\delta(x-x'),
\end{equation}
求 $F_{p,p'}$。
\end{problem}

\begin{answer}
% 在此处作答
\end{answer}

\end{document}
